(1) 理想无损双输出干涉仪的两个计数端口在给定相位差 \(\phi\) 下探测到独立的相干态泊松光子数 \(N_+,N_-\)。一次测量周期的总平均探测光子数为 \(N\)，两个端口的平均数为 \(\mu_\pm=N(1\pm\cos\phi)/2\)。由泊松分布独立性求差分 \(D=N_+-N_-\) 的均值、方差，解释强本振不意味着噪声能消失。
(2) 在非奇异相位 \(\phi\notin\pi\mathbb Z\) 处，由完整的联合计数似然 \(P(N_+,N_-|\phi)\) 计算经典Fisher信息和Cramér--Rao相位标准差下界；比较直接线性化差分的噪声误差传播，说明端点的局部不可辨识性。
(3) 设对比度 \(0<V\le1\) 使 \(\mu_\pm=N(1\pm V\cos\phi)/2\)，总平均计数仍为 \(N\)。求Fisher信息，证明 \(\phi=\pi/2\) 是 \(V<1\) 的最优局部工作点并给相位灵敏度；若总入射光子数为 \(N_{\rm in}\)、均匀探测效率为 \(\eta\)，替换为实际 \(N\) 并说明损失不能凭借重归一化消除。
(4) 写出使上述估计接近极限的高计数局部估计器，给出偏置和动态范围的限制；对 \(\phi=\pi/2\) 核验用差分测量得到的信噪比及维度，并澄清经典泊松计数极限与压缩光/NOON态不同。
